\section{Motivation: recovery and decision-rule validation}
\label{sec:intro}

A reporting rule can withhold unreliable estimates only if its thresholds
are calibrated to the uncertainty of the reported quantity. We examine a
50-cell rule in colorectal single-cell analysis through three empirical tests:
coverage against a fixed empirical source-pool target, rejection under the
null, and stability of first-crossing candidates. The study is a self-audit
of a previously overstated validation claim in an unpublished project. Its
contribution is the measured behavior of this rule and the reproducible target
comparison, rather than a new principle connecting point and interval accuracy.

Our motivating application separates two explanations for differentiation-marker
loss in colorectal cancer: fewer marker-producing cells (a \emph{compositional}
change), or lower output from the cells that remain (an \emph{intrinsic}
change). These explanations suggest different biological interpretations.
An empirical comparison of per-cell means requires the relevant population
to be observed in both tissues. Let $f_N,f_T$ denote the relevant cell fractions
in reference and diseased tissue, and $m_N,m_T$ their mean per-cell output.
A reference-based decomposition, related to demographic standardisation and
Oaxaca--Blinder decompositions \citep{kitagawa1955,oaxaca1973,blinder1973}, is
\begin{equation}
\underbrace{(f_T-f_N)m_N}_{\text{compositional}} +
\underbrace{f_N(m_T-m_N)}_{\text{intrinsic}} +
\underbrace{(f_T-f_N)(m_T-m_N)}_{\text{interaction}}
= f_Tm_T-f_Nm_N.
\label{eq:kitagawa}
\end{equation}
The interaction is reported separately. If no relevant cells are observed in
an arm, its sample mean is undefined; the identity cannot supply the missing
mean. In a synthetic population with no such cells, the corresponding
conditional mean itself has no referent. Substituting zero in either case
turns an unavailable contrast into an apparent effect.

The implemented reporting rule returns an estimate at $n\geq50$, flags a wide
interval at $20\leq n<50$, and returns \notest{} below 20, where $n$ is the
relevant-cell count in diseased tissue. Unavailable reported contrasts are
stored as \texttt{None}, not zero. The diagnostic oracle calculations examined
below also retain raw arithmetic before this reporting gate. We evaluate the boundary directly and retain pre-gate intervals at smaller
counts to diagnose the curve's behavior.

\paragraph{Relation to existing validation methods.}
Simulation-study methodology already separates estimands, methods and
performance measures \citep{burton2006,morris2019}. Generator--analysis alignment
is discussed in congeniality \citep{meng1994}, plasmode and semi-synthetic
benchmarks \citep{franklin2014,curth2021}, interventional evaluation
\citep{gentzel2019} and the inverse crime \citep{wirgin2004,kaipio2007}; generators
can also disadvantage good estimators \citep{shaw2026}. We apply these distinctions to the recorded reporting-rule decision. For clinical trials, external-control validation
already examines bias and type-I error using held-out studies
\citep{ventz2019}. Prognostic adjustment with historical data retains trial
randomisation and evaluates precision gains \citep{schuler2022}; target-trial
emulation states the causal question and design explicitly \citep{hernan2016}.
These methods motivate the clinical-simulation validation questions in
Section~\ref{sec:validation}. The empirical study uses single-cell pools; a separate synthetic
external-control experiment illustrates target-population mismatch. Neither
validates a patient world model. Simulation-based calibration evaluates posterior ranks
\citep{talts2018,modrak2025}.
Appendix~\ref{app:relatedwork}
provides further context.
